
\section{Ethical Statement}

\benchmark aims to drive embodied AI solutions that fulfill human need. A primary ethical consideration of \benchmark is therefore the method used to determine such need. To understand which activities would be useful for humans, we survey 1,461 respondents on Amazon Mechanical Turk and collected 50 responses for each activity, aiming for a clear consensus signal. Though this is a large group, the results are still biased by the fact that the researchers, annotators, and data providers only represent a small population relative to potential users of such technology. We plan to address these biases by making \dataset open-source and invite a wider community to contribute to its knowledge base. 

Furthermore, compared to the U.S. population, the demographic representation in our survey skews white, male, non-disability status, mid-five figure incomes, and 30-40 age range with more representation on the higher side than the lower side of that range\textemdash closer to the demographics of Mechanical Turk workers. This also biases the survey population toward certain responses, creating potential ethical limitations. 

Specifically this bias may affect the survey's representation of people who would be affected most by autonomous agents~\cite{mckinsey2019automation}. We therefore asked several explicit questions such as ``Do you do household work for a living?'', ``If you do household work for a living, would you benefit from assistance?'',  ``Do you pay for someone else to do household work for you?''. Future work involves using these questions to direct benchmark development. 


